\section{Construction details and intermediate counts}

This appendix brings together the construction and proof explanations that
accompany the short note. The Lean definitions and theorem statements determine
the exact hypotheses and conclusions. The explanations below do not constitute
a human inspection of the formalization.

\subsection{Four-state and block constructions}

Let an original machine have $m$ states and $n$ symbols. Write $m_+$ and $m_-$
for the numbers of states entered by right and left moves, respectively.
Morita's four-state conversion uses $(2m+1)n$ symbols. Giving the two directions
separate number ranges reduces the expression to $(m+1)n$.
Counting in reverse while carrying a symbol to the left gives the upper bound
$(\max(m_+,m_-)+1)n$.
Splitting states into blocks, merging roles, and removing unused
symbol--counter pairs gives $n+\sum_{i\in V}|X_i|$, where $V$ is the set of
counter values used and $X_i$ is the set of original symbols used with value $i$.
For the conversion from $(10,8)$, the respective symbol counts are 168, 88,
48, and 46. The last value is the block construction before the later
symbol-merging improvements.

The corresponding formal statements are collected in
\texttt{LeanUniv.FinalFormulas}. The expression for reverse counting is an
upper bound; the other formulas just listed are equalities for the stated
conversion functions. These formulas concern construction families and do
not give lower bounds over all universal machines.

For the final five symbol-merging points, the original machine is $(13,7)$.
The recipe files specify the blocks and counter choices, and the merge files
specify classes of symbols. The generators reconstruct the tables, while
Lean checks their finite properties and applies the general symbol-merging
theorem. Rebuilding these files reproduces the selected constructions; it
does not repeat the optimization search or establish optimality.

\subsection{Three-state conversion}

For Morita's three-state conversion, let
$M=m_+n+2n+2$ and $N=mn^2+n^2+1$ be the sizes of the intermediate machine.
The original symbol-count formula is $N(2M-1)$. Sharing the up and down marks
gives $NM$. A further construction assigns a separate counter limit to each
symbol; its count is $N$ plus the sum of these limits.

For the source $(32,3)$, the successive recorded symbol counts are 36,654,
18,476, 1,820, 659, and 600. The 1,820-symbol intermediate table has only a
size calculation: adequacy of its counter limits throughout every run was
not proved. The 659-symbol table rebuilds the intermediate machine using
one symbol per cell. Reflecting the source machine left to right gives the
600-symbol table. Simulation and universality of the final $(3,600)$ witness
are established by the corresponding Lean theorems.

\subsection{Two-symbol conversion and automated search}

For a conversion using $k$-bit codes, let $P_r$ and $P_w$ be the sums, over
states, of the numbers of read-code and write-code prefixes of lengths
$1,\ldots,k-1$. Let $D_0$ be the set of source states having a rule whose
target is also a source state. The formalized state-count formula is
$|D_0|+(k-1)m+P_r+P_w$, bounded above by $mk+P_r+P_w$.

From $(13,7)$, the recorded state counts are 169 using Morita's code,
149 with a search-selected code, 135 after selecting written values and path
types, 92 after state merging, and 85 after the broader automated search.
The solver that selected the intermediate tables is not itself a verified
Lean conversion function. Its concrete output tables are checked in Lean;
the proof for the merged table also establishes the required relationship
to its pre-merge table. The relevant intermediate-value theorems are in
\texttt{LeanUniv.FinalValues}.

The automated search was modeled on Karpathy's autoresearch
(\url{https://github.com/karpathy/autoresearch}).
An interactive AI session prepared the search program, fixed evaluator,
and instructions at the human's request. After the human started a session,
the agent edited the search program and the evaluator checked candidates
for determinism, reversibility, and finite CTSH simulations. Only candidates
subsequently proved in Lean counted as results. The search found $(6,37)$,
$(8,28)$, $(9,27)$, and $(85,2)$; only the last remains non-dominated after
the symbol-merging improvements. The repository includes the construction
recipes needed to rebuild the final witnesses.

\section{Outline of the two-state decidability proof}

The formal entry points are \texttt{halting\_decidable\_uniform} and
\texttt{no\_le\_two\_state\_URTM}. This section summarizes the structure
of the existing proof; the formal statements specify its scope.

First remove transitions into states with no outgoing transitions by the
reduction described in the main text. Let $[a,b]$ contain the initial head
position and all initially non-blank cells, and put $w=b-a+1$.
Let $\Sigma$ contain the blank and every symbol read or written by the
table, and let $n=|\Sigma|$. Consider an initial ID whose tape symbols
lie in $\Sigma$.

\begin{enumerate}
\item In a reversible table, transitions entering the same state move in
the same direction. If a two-state table has both left and right moves,
one state is entered only from the left and the other only from the right;
there is no remaining state for a zero-shift target. The state therefore
records the direction of the last move. The transition function can be
represented by a partial injection on direction--symbol pairs.

\item Consider a half-line to the right of a fixed boundary. Every entry
has the right-moving state and every exit the left-moving state. While the
head is outside the half-line, its contents do not change. The outcome of
the next entry therefore depends on its initial contents and on the number
of earlier entries, rather than on the outside tape. The only information
passed back to the outside is the return itself.

\item For initially blank nested half-lines, their local computations are
translations of one another. In a halting computation, each visited deeper
half-line has fewer entries than its predecessor. Otherwise the computation
repeats translated to the right and never halts. This bounds the maximum
depth by the number of entries at the outer boundary.

\item The sequence of visits to a cell follows a finite local clock on
direction--symbol pairs. There are $2n$ such pairs. This bounds by $2n+1$
the number of entries into a blank half-line if the computation halts
inside it. The left-hand case follows by reflection.

\item Replace excursions outside $[a,b]$ by immediate returns at its ends.
The resulting dynamics within the window is finite. In the existing proof,
$Z=n^w\cdot 3w$ bounds the window states, and
$N=Z(2n+3)+1$ bounds the relevant entries into each outside half-line.
Combining this with the depth bound confines a halting run to
$[a-N,b+N]$.

\item That interval contains finitely many IDs over the finite alphabet.
A halting run of a deterministic machine cannot repeat an ID. Thus a
computable bound on the number of steps suffices to decide halting by
bounded simulation. The bound need not be practical.
\end{enumerate}

If only one movement direction occurs, once the head leaves the initial
window it repeatedly reads blanks. Repetition of a state gives a translated
repetition of the run. A halting run can therefore be bounded in this case
as well. The general nonexistence theorem combines the computability of
halting for these tables with Mathlib's undecidability of the halting problem.

The first step is where reversibility restricts the two-state model. With
three states, a half-line can receive different kinds of entries moving in
the same direction; the two-state argument does not apply unchanged.

\section{Outline of the three-state two-symbol proof}

The proof normalizes the tables and checks the 359,829 keys of the normal
forms with the Lean kernel. It handles most tables by certificates constraining visits
outside a window and by finite-state bounds. Other cases reduce to bounded
halting for a two-state machine after removing a state, or require a
separate argument about their behavior.

Counts of equivalence classes up to renaming and left--right reflection
were computed outside Lean: there are 26,861 such classes in the recorded
enumeration, and the last ten classes contain 120 tables. The AI supplied
arguments for these final classes, proving that a halting run stays within
a bounded neighborhood of the initial window.

The resulting theorem is
\texttt{URTM.Meta.no\_three\_two\_URTM\_all}, followed by the transfer to
\texttt{no\_three\_two\_URTMStrict\_all}. The original target condition is
the string-ID form of universality with the stated restriction on encoded
symbols. The transfer from \texttt{IsUTMStrict} is essential: a proof of
nonexistence for a different universality definition alone would not imply
the claim made in the note. The finite enumeration files and all imported
proof dependencies are included in the repository.

\section{Comparison with prior work and verification scope}

Morita's nine machines provide the baseline~\cite{Morita17}. Six of these
points remain on the frontier together with the seven new points. The work
improves conversion and symbol-assignment choices; it does not prove that
the resulting state--symbol pairs are globally smallest.
Kudlek and Rogozhin's merging of a blank and another symbol in an ordinary
UTM~\cite{KudlekRogozhin02} guided the symbol-merging stage. The general
theorem here makes the preservation conditions explicit for the formalized
universality predicate, including recovery of the symbols on halting tapes.

The lower-bound region quoted for standard UTMs comes from the survey by
Woods and Neary~\cite{WoodsNeary09}. Its conventions for halting and
zero-shifts differ from the present model. The new nonexistence claims are
therefore established directly in the stated reversible model rather than
inferred solely from that comparison. The AI also used prior work on
irreversible machines to guide the three-state proof; the author has not
checked that correspondence against the original sources.

The supplied generators read Morita's original test configurations from
the machine descriptions attached to his 2015 slides. For eight machines,
Lean reproduces 44 published halting runs, with the count 1,073 in one
$(13,7)$ comment corrected to 1,173. For $(3,36654)$, the exact step counts
are 471,798, 2,348,492, and 57,220. The large table is checked through a
lookup function proved to agree with it and certificates for successive
run segments. These finite examples support comparison of the operational
model with the supplied data, while the general simulation and universality
theorems cover arbitrary encoded inputs.

The two proof environments reuse the specification by copying it.
The repository checks that the copies of \texttt{Spec.lean} are byte-identical
and that the copies of \texttt{Strict.lean} differ only in their imports.
Each environment is built separately and its imported declarations are
audited for nonstandard axioms. These checks establish properties of the
formal artifacts; they do not establish the faithfulness of the formalization
to the intended mathematical claims.
